% Generated draft: requires mathematical and source review.
\documentclass[11pt]{article}
\usepackage[margin=25mm]{geometry}
\usepackage{fontspec}
\setmainfont{Latin Modern Roman}
\usepackage{amsmath,amssymb,mathtools,unicode-math}
\setmathfont{Latin Modern Math}
\usepackage{enumitem,needspace,xurl,hyperref}
\setlength{\parindent}{0pt}
\setlength{\parskip}{5pt}
\setlength{\emergencystretch}{3em}
\newcommand{\N}{\mathbb{N}}
\newcommand{\Z}{\mathbb{Z}}
\newcommand{\Q}{\mathbb{Q}}
\newcommand{\R}{\mathbb{R}}
\newcommand{\C}{\mathbb{C}}
\newcommand{\F}{\mathbb{F}}
\newcommand{\A}{\mathbb{A}}
\newcommand{\PP}{\mathbb{P}}
\newcommand{\E}{\mathbb{E}}
\newcommand{\eps}{\varepsilon}
\newcommand{\dd}{\,\mathrm{d}}
\begin{document}
\section*{UnsolvedMath ID: 11000055}
\textbf{Source problem code:} catalogue entry
\subsection{Definitions and mathematical statement}
For \(g\ge2\), let \(\operatorname{Mod}_g\) be the source's mapping class group, conventionally that of a closed oriented genus-\(g\) surface. An involution is a nonidentity element \(t\) with \(t^2=1\). Define
\[
 \iota(g)=\min\left\{m:\operatorname{Mod}_g=\langle t_1,\ldots,t_m\rangle,
                  \ t_j^2=1,\ t_j\ne1\right\},
\]
with \(\iota(g)=\infty\) if no such finite set exists. This convention does not assume generation by involutions in every genus.

Problem 5.8 asks for sharp bounds on \(\iota(g)\) for each \(g\ge2\), including matching constructions and obstructions sufficient to establish sharpness. In particular it asks, for each \(g\ge7\), whether there exist three involutions generating \(\operatorname{Mod}_g\), equivalently whether
\[
 \iota(g)\le3.
\]
This is ordinary subgroup generation, not normal generation, and it is not a request for a single uniform triple acting in different genera. The surface/boundary conventions must be confirmed, since they affect torsion. The supplied record gives no numerical conjectural answer or previously established bounds, so none is inserted.
\subsection{Short English statement}
Find the sharp minimum number of involutions generating the mapping class group in every genus at least two. In particular, decide whether three suffice in each genus at least seven. Generation by involutions is not presumed when defining the minimum.
\subsection{Sources}
\begin{itemize}
\item Local UnsolvedMath registry, permanent ID \(11000055\), exact source code \texttt{catalogue entry}, Problem 5.8, ``Sharp bounds for involution generating sets,'' PDF page 45. No individual proposer or proposed year is supplied.
\item Farb, edited volume \emph{Problems on Mapping Class Groups and Related Topics}: \url{https://www.math.uchicago.edu/~farb/papers/mcgbook.pdf}. The volume contains 329 explicitly typeset numbered items across contributed chapters; repeated numbering is distinguished by PDF page.
\item Extraction pdf-font-structure, accessed 2026-07-29. Source evidence NEEDS\_REVIEW presents an open problem and requires a current primary-literature audit. Registry status open and classification OPEN-TRIAGE; the summary generically mentions literature bounds and possibly open sharp versions, but supplies no references or exact numbers and says ``Without statement, mark open.'' This limited assessment is not a verified current status determination of the fully supplied question.
\item Dataset attribution: Ulam AI / UnsolvedMath Contributors, CC BY 4.0. Source rights NEEDS\_REVIEW: publicly accessible author/publisher copy, redistribution terms requiring review.
\end{itemize}
\end{document}
